(1) 令 \(\gamma=1/\sqrt{1-v^2/c^2}\)。两粒子四动量分别为 \(p_1^\mu=(\gamma mc,\gamma mv)\) 和 \(p_2^\mu=(\gamma mc,-\gamma mv)\)，其和为 \(P^\mu=(2\gamma mc,0)\)；总能量 \(E_{\rm tot}=2\gamma mc^2\)，总动量为零。

(2) 因碰撞中四动量守恒，末态物体必须静止于这个实验室系，即 V=0。不变质量由 \(M^2c^4=E_{\rm tot}^2-c^2P_x^2\) 给出
\[
M=2\gamma m.
\]
非弹性碰撞中的复合静质量大于两粒子静质量之和，这是储存内部能量的结果。

(3) 与两颗分离的静质量之和比较，末态额外的内部能量为
\[
\Delta E_{\rm int}=(M-2m)c^2=2mc^2(\gamma-1).
\]
碰前总动能正好为 \(K_{\rm ini}=2(\gamma-1)mc^2\)，所以全部相对运动动能转成了复合物内部能，既没有丢失能量也没有必须对外辐射的能量。

(4) 对 \(v/c\ll1\)，\(\gamma=1+v^2/(2c^2)+O(v^4/c^4)\)，故 \(\Delta E_{\rm int}=mv^2+O(v^4/c^2)\)，恰为两颗牛顿粒子动能总和。\(v\to c^{-}\) 时 \(\gamma\to\infty\)、M 无界，显示无限增大动能会使复合系统不变质量无限增大。
